\paragraph{The scalar auxiliary section before the definite lift.}
\label{par:selected-auxiliary-cusp}

The definite last lift is treated separately. Before reaching it one needs an auxiliary scalar section with controlled values for the height calculation.

Let \(\Lambda\) be the rank-\(r\) polarized lattice of the fixed lift. For this construction use the odd integer \(R=5\). Write
\(\Gamma_R=\operatorname{Aut}_{\mathcal O_E}(\Lambda/R^4\Lambda)\).
The module modulo \(R^4\) is free of rank \(r\) over
\(\mathcal O_E/R^4\mathcal O_E\), since \(\Lambda\) is projective of
constant rank and the quotient is a product of Artinian local
rings. The kernel of reduction of the matrix group modulo \(R\)
has order \(R^{6fr^2}\). Its elements are all matrices of the
form \(1+A\) with entries of \(A\) in the nilpotent ideal
\((R)/(R^4)\). The reduction contains \(-1\) of order two.

Consequently 
\[
  |\Gamma_R|=2R^4 i_\eta
\]
 for an integer
\(i_\eta>0\). Also \(|\Gamma_R|\le R^{8fr^2}\) by counting all
matrices. Choose lifts \(x_b\), \(1\le b\le r\), of a basis
modulo \(R\). Index tuples by
\(J=(\Lambda/R\Lambda)^r\). Choose representatives of their
entries and put
\[
             x_{l,b}=x_b+R\widetilde l_b\pmod {R^4\Lambda}
             \qquad(l\in J).
\]
For each \(\gamma\in\Gamma_R\), the elements \(\gamma x_{l,b}\)
span the full module modulo \(R^4\): they generate modulo the
nilpotent ideal \((R)\), then generate by Nakayama's lemma.
There are \(R^{2fr^2}\) tuples.
\label{cor:selected-auxiliary-cusp-1}

\begin{auxlemma}[A nonzero theta value on a torsion translate]
\label{lem:eta-torsion-translate}
Let \(A/\mathbb C\) be a principally polarized abelian variety, let
\(L\) represent its polarization and let \(s\) generate \(H^0(A,L)\).
For every integer \(R\ge2\), the torsion translates have empty common zero locus:
\[
 \forall x\in A(\mathbb C)\quad\exists t\in A[R]\quad s(x+t)\ne0 .
\]
\end{auxlemma}
\begin{proof}
Write \(p=[R]\). Its deck group is \(A[R]\), so the character
decomposition of the finite étale cover and the projection formula give
\[
 H^0(A,p^*L)=
 \bigoplus_{P\in\operatorname{Pic}^0(A)[R]} H^0(A,L\otimes P).
\]
Indeed the character sheaves of \(p_*\mathcal O_A\) are the line
bundles trivialized by \(p\), each with multiplicity one. The
principal polarization identifies the indexing sets by
\[
 A[R]\xrightarrow{\ \sim\ }\operatorname{Pic}^0(A)[R],
 \qquad t\longmapsto T_t^*L\otimes L^{-1}.
\]
Thus each summand is spanned by the pullback of a translate of \(s\).

The line \(p^*L\) is generated by its sections. To see this, the
theorem of the square gives \(p^*L\simeq L^{R^2}\otimes P_0\) for
some \(P_0\in\operatorname{Pic}^0(A)\). At any specified point take
a product of \(R^2\) translates of \(s\), with the sum of the
translation points fixed by \(P_0\). Each individual vanishing
condition cuts out a proper closed subset of that fiber of the
addition map, which is isomorphic to \(A^{R^2-1}\). A product
avoiding them all has the required line bundle and a nonzero value.
Consequently at a lift of \(x\) some character summand is nonzero.
This gives the claimed translate.
\end{proof}

\begin{auxlemma}[Decay along a shortest CM period]
\label{lem:eta-period-decay}
Let \(V/L\) be a complex torus of dimension \(g=fr\) with its
principal Riemann form \(B=\operatorname{Im}H\), Euclidean norm
\(\|v\|^2=H(v,v)\), and polarized \(\mathcal O_E\)-action with
Rosati involution given by conjugation. Let
\(u\in\mathcal O_E\) satisfy \(u+\bar u=1\). Use the theta line
whose semicharacter is \(k\mapsto(-1)^{B(k,uk)}\), and let
\(s\) be a section of Haar \(L^2\)-norm one. Write
\(\lambda=\min_{0\ne k\in L}\|k\|\),
\(\rho=|\operatorname{disc}E|^{1/(2f)}\) and
\(\Omega=2f^{3/2}\rho\). Suppose that
\[
 M\ge\rho^2,\quad \rho\ge16,\quad
 M^{-1/2}\le h\le 1/\Omega,\qquad
 P=\begin{cases}
 2f/(\rho^2\lambda^2),&\lambda<h,\\
 0,&\lambda\ge h.
 \end{cases}
\]
Let \(R\ge2\), and let \(a_1,\ldots,a_r\in V\) satisfy
\(R^2a_b\in L\) with these periods generating \(L/RL\) over
\(\mathcal O_E\). Put
\[
 C_0=r\bigl[3+g\log(2\sqrt{gM})
                  +2f\log(2\sqrt{2f})\bigr].
\]
For the pointwise theta metric,
\begin{equation}\label{eq:eta-period-decay}
 \prod_{b=1}^r\|s(a_b)\|
 \le e^{C_0}\max(1,P)^{rg/2}
                 \exp\!\left(-\frac{\pi P}{2R^4}\right).
\end{equation}
\end{auxlemma}
\begin{proof}
\proofstep{A bound valid on the whole torus.}
Put \(S=\sup_z\|s(z)\|\) and \(b_0=\min(1,\lambda/2)\).
A polydisc of coordinate radius \(b_0/\sqrt g\) in orthonormal
coordinates injects into the torus. Center the Appell--Humbert
frame at the point in question. On this polydisc the squared
metric is \(e^{-\pi\|z\|^2}\) times the squared absolute value.
The principal alternating form makes the period covolume one in
these coordinates, so Euclidean measure on a fundamental domain
is Haar probability measure. Submean on the polydisc gives
\[
 S\le e^{\pi/2}\left(\frac{\sqrt g}{\sqrt\pi\,b_0}\right)^g
 \le e^2 (2\sqrt{gM})^g\max(1,P)^{g/2}.
\]
In fact \(\lambda^{-1}\le\sqrt M\) when \(P=0\), and
\(\lambda^{-1}=\rho\sqrt{P/(2f)}\) otherwise. This already proves
\eqref{eq:eta-period-decay} for \(P=0\).

\proofstep{An isotropic subspace and its frequencies.}
Choose a shortest period \(\xi\) when \(\lambda<h\). There are
integers \(\omega_j\in\mathcal O_E\) spanning \(E\) over
\(\mathbb Q\) with \(\max_\sigma|\sigma\omega_j|\le\Omega\).
Indeed, equip \(E\otimes\mathbb R\) with the squared norm
\(\sum_\sigma|\sigma x|^2\). The dual of \(\mathcal O_E\)
is its codifferent for this metric. Every nonzero member of the
codifferent satisfies
\[
 |N_{E/\mathbb Q}y|\ge |\operatorname{disc}E|^{-1},
 \qquad \sum_\sigma|\sigma y|^2
       \ge 2f |N_{E/\mathbb Q}y|^{1/f}\ge 2f/\rho^2 .
\]
Thus the dual first minimum is at least \(\sqrt{2f}/\rho\).
Euclidean transference gives last minimum of \(\mathcal O_E\)
at most \(2f\) divided by the dual first minimum
\cite{banaszczyk}. This is at most \(\Omega\).

For each \(j\), integrality of the alternating form and the
bound on the operator norm of \(\omega_j\) give
\[
 |B(\xi,\omega_j\xi)|\le\Omega\lambda^2<1,
 \qquad B(\xi,\omega_j\xi)=0.
\]
By linearity and the Rosati identity, \(W=\operatorname{span}_{\mathbb R}
(\mathcal O_E\xi)\) is isotropic. Set \(I=W\cap L\) and
\(n=\dim_{\mathbb R}W\le2f\). The semicharacter is one on \(I\),
since \(uI\subset I\). With \(I^*\) the Euclidean dual in \(W\),
every nonzero \(p\in I^*\) has nonzero integral pairing with some
\(\omega_j\xi\). Hence
\[
 \|p\|\ge(\Omega\lambda)^{-1}\ge1.
\]

\proofstep{The Fourier expansion on the isotropic subspace.}
For \(a\in V\) let \(U(a)\in W\) represent
\(k\mapsto B(a,k)\) by the Euclidean inner product. In the
Appell--Humbert frame centered at \(a\), restrict the section to
\(W+iW\). Multiply its holomorphic representative by
\(\exp(-\pi(z,z)/2)\), where \((\ ,\ )\) is the complex bilinear
extension of the real inner product on \(W\).
The resulting function \(G_a\) satisfies
\[
 |G_a(0)|=\|s(a)\|,\qquad
 |G_a(x+iy)|\le S e^{\pi\|y\|^2}\qquad(x,y\in W).
\]

Its translation multiplier on \(I\) is \(e^{\pm2\pi iB(a,k)}\),
with the sign fixed by the convention for \(H\). For a frequency
\(p\in I^*\pm U(a)\), move the Fourier coefficient integral from
\(W/I\) to \((W-ip)/I\). Opposite faces cancel by the multiplier
law, and the exponential from the Fourier character gives
\[
 |\widehat G_a(p)|\le S e^{\pi\|p\|^2-2\pi\|p\|^2}
                    = S e^{-\pi\|p\|^2}.
\]

These bounds give absolute convergence at the origin. Reversing all
frequencies if necessary, we may therefore sum over \(U(a)+I^*\).

\proofstep{A norm bound for every tuple of frequencies.}
Take arbitrary \(p_b\in U(a_b)+I^*\). Since \(I\) is saturated in
\(L\), every integer functional on \(I\) extends to \(L\). By the
principal polarization \(U(L)=I^*\), so we can write
\[
 R^2p_b=U(k_b),\qquad k_b=R^2a_b+R^2t_b,\quad t_b\in L .
\]
The \(k_b\) still generate \(L/RL\). The integral functional
\(\alpha\mapsto B(k_b,\alpha\xi)\) on \(\mathcal O_E\) has the
form \(\operatorname{Tr}_{E/\mathbb Q}(y_b\alpha)\) for an element
\(y_b\) of the codifferent. At least one \(y_b\) is nonzero.
Otherwise, by the Rosati identity, every element of
\(\sum_b\mathcal O_E k_b+RL=L\) would pair with \(\xi\) in
\(R\mathbb Z\). Perfectness excludes this because a shortest
period is primitive.

The norm of this trace functional for the embedding metric is
\((\sum_\sigma|\sigma y_b|^2)^{1/2}\). Its value at \(\alpha\) is
\((R^2p_b,\alpha\xi)\), and
\(\|\alpha\xi\|\le\|\xi\|(\sum_\sigma|\sigma\alpha|^2)^{1/2}\).
For a nonzero \(y_b\) the preceding codifferent estimate gives
\[
 R^4\|p_b\|^2\lambda^2\ge
 \sum_\sigma|\sigma y_b|^2\ge 2f/\rho^2,
 \qquad \sum_b\|p_b\|^2\ge P/R^4 .
\]

Reserve half of each Gaussian exponent for this last inequality.
The rest is summable uniformly on every translate \(D\) of \(I^*\).
Indeed the balls of radius \(1/2\) about \(p\in D\) have disjoint
interiors. On such a ball \(\|p\|^2\ge\|x\|^2/2-1/4\), whence
\[
 \sum_{p\in D} e^{-\pi\|p\|^2/2}
 \le \frac{e^{\pi/8}}{\operatorname{vol}B_n(1/2)}
                  \int_{\mathbb R^n}e^{-\pi\|x\|^2/4}\,dx
 \le e^{\pi/8}(2\sqrt n)^n .
\]

The last inequality uses the inscribed cube of side \(1/\sqrt n\).
Multiplying the absolutely convergent Fourier bounds for all
\(b\), then using \(n\le2f\) and the bound for \(S\), proves
\eqref{eq:eta-period-decay}.
\end{proof}

\begin{auxlemma}[Parameter bounds for the two theta estimates]
\label{lem:joint-theta-parameters}
Let \(g=fr\), \(D=D_F\) be the differential degree of the light weight,
\(M=g_0g_1\), and \(\rho=\operatorname{rd}(E)\). Suppose
\[
 \begin{gathered}
 f,r,m,D\ge1,\quad g\ge2,\quad m^4\le f,\quad \rho\ge\max(16,f^2),\\
 R_0^2=90(g+D),\qquad M\ge R_0^2\rho^{28}Dr .
 \end{gathered}
\]
With \(d=2fm\), define
\[
 \Omega=2f^{3/2}\rho,\qquad h^2=fR_0^2\rho^{21}/M,\qquad
 B=M/(R_0^2\rho^{28}),\qquad T_0=MB .
\]
Then \(M^{-1/2}\le h\le\Omega^{-1}\), and all six estimates used
in the lower bound for the light determinant hold simultaneously:
\begin{align}
 d^2\Omega h^2&<1, & B&\ge1,\label{eq:joint-theta-smallness}\\
 90m^2d^2\rho&\le M/2, & D&\le M,\label{eq:joint-theta-dimension}\\
 4D(7+\log(MD))+d+1 &\le T_0,\label{eq:joint-theta-log-cost}\\
 g\log2+2d+3+(D+1)(r\log r+3r+3)&\le T_0.\label{eq:joint-theta-rank-cost}
\end{align}
\end{auxlemma}
\begin{proof}
We have \(d\le2\rho\), \(m^2\le\rho\), and
\(f^{3/2}\le f^2\le\rho\). The common level bound gives
\[
 B\ge Dr,\qquad T_0\ge MDr,\qquad R_0^2\ge270.
\]
In particular \(D\le M\) and \(h\sqrt{M}\ge1\). The two
smallness requirements use different powers of \(\Omega\):
\[
 d^2\Omega h^2\le\frac{8R_0^2\rho^{26}}M\le\frac8{\rho^2}<1,
 \qquad (\Omega h)^2\le\frac{4 f^4}{\rho^5Dr}\le1.
\]
Also \(90m^2d^2\rho\le360\rho^4\le M/2\).

For the logarithmic allowance use \(D\le M\) and
\(\log M\le M/32\) for \(M\ge256\), by monotonicity of
\(t/32-\log t\) in that range. Since \(M\ge270\rho^{28}\),
\[
 4D(7+\log(MD))+d+1
 \le D(29+2\rho+M/4)\le MD\le T_0 .
\]
Finally \(r\le g\), \(\log r\le g\), and \(D+1\le2D\) give
\[
 g\log2+2d+3+(D+1)(r\log r+3r+3)
 \le (3g+4\rho+15)Dr\le MDr\le T_0 ,
\]
using \(M\ge90(g+D)\rho^{28}Dr\) in the second inequality.
\end{proof}

\begin{proposition}[The scalar cusp section before the definite lift]
\label{prop:auxiliary-scalar-section}
For each nondefinite lift the scalar construction gives a cusp section
\(\eta\) nonzero on every component with the same finite and
archimedean boundary comparisons used for the light columns. For the
page parameters, one may take the decay constant \(P_2\) and the upper bound \(P_\eta\) for the mean log supremum divided by \(i_\eta\) to satisfy
\[
 \max\{P_2,P_\eta\}\le200000(frL_0)^2,
\]
where \(L_0\) is the run's logarithmic scale.
\end{proposition}

\begin{proof}
\proofstep{The scalar section before descent.}
Over the characteristic-zero torsor of full linear level \(R^4\), denote the
normalized odd-level theta null values by \(A(t)\) for
\(t\in\Lambda/R^4\Lambda\). They are the values without differentiation at the point \(t\) of order dividing \((R^2)^2\)
with the odd-level normalization of order \(R^2\). Use
\begin{equation}\label{eq:selected-eta-product}
 \eta=\prod_{\gamma\in\Gamma_R}
             \left(\sum_{l\in J}m_l
                       \prod_{b=1}^{r}A(\gamma x_{l,b})\right).
\end{equation}
The frame law gives each theta null value central degree one
and zero row degrees. The torsor action permutes the outer
factors, so their product descends to the original coefficient
field and level. Its weight is 
\[
  r|\Gamma_R|u_0=i_\eta g_0^*u_0,
\]
with \(g_0^*=2\cdot5^4r\).
\label{cor:selected-auxiliary-cusp-2}

Choose one frame at the chosen complex embedding on each
component. For a torsion label \(t\), rigidification gives a nonzero frame
at \([R^2]t\), hence
\[
 A(t)\ne0\quad\Longleftrightarrow\quad s([R^2]t)\ne0.
\] At each tested frame and each \(\gamma\),
the shifts in a single column give the following points:
\[
  R^2(R\widetilde l_b)=R^3\widetilde l_b,\qquad R(R^3\widetilde l_b)=0\pmod{R^4}.
\]
As \(\widetilde l_b\) ranges through representatives modulo \(R\),
these last classes cover all of \(A[R]\), including after the
automorphism \(\gamma\). Applying Lemma~\ref{lem:eta-torsion-translate}
at \([R^2]\gamma x_b\) gives a nonzero value for that column. The shifts
of the columns are independent and all occur in \(J\); hence the
linear functional on \((m_l)_l\) in that factor is nonzero.

With \(M_c\) tested frames, the box \(0\le m_l\le M_c|\Gamma_R|\)
contains an integral point on which every such functional is
nonzero. Set \(B=M_c|\Gamma_R|\). A proper hyperplane contains at most \((B+1)^{|J|-1}\) points of the box, so the number left after excluding all \(B\) hyperplanes is at least
\[
  (B+1)^{|J|}-B(B+1)^{|J|-1}=(B+1)^{|J|-1}>0.
\]
 The odd-level test requires
characteristic zero, the full \(R^4\) torsion frame and the
rigidified theta line. All occur in the chosen complex uniformization. For the chosen test point \(x_a\) on component \(a\) the coefficients give
\[
 \prod_{\gamma\in\Gamma_R}
   \left(\sum_l m_l\prod_b A(\gamma x_{l,b})(x_a)\right)\ne0.
\]
Thus \(\eta\) is nonzero on every component. They are used at every finite and infinite
place.
\label{cor:selected-auxiliary-cusp-3}

\proofstep{Finite orders in an integral frame.}
At a finite algebraic test point, pass to a split semiabelian
uniformization in a finite valued extension. Write \(Q\) for
the nonnegative quadratic form of the integral theta expansion in
the degeneration lattice and \(y_t\) for the lattice class of a torsion evaluation at level \(R^4\). In a frame with integral theta
generator, we claim that
\[
       \operatorname{ord} A(t)
       \ \ge\ \frac{1}{R^4}
                   \min_z Q(y_t+R^2z).
\]

To compute the order, let \(m=R^2\), and put
\(B_Q(y,z)=Q(y+z)-Q(y)-Q(z)\). Use the period and character
lattices identified by the principal polarization. In the symmetric
Mumford degeneration data, a period \(y\) contributes
\(B_Q(y,z)\) to the order of character \(z\) and \(Q(y)\) to
the order of its theta multiplier. Indeed the integrally normalized Poincar\'e
biextension has pairing of order \(B_Q(y,z)\). Pullback along
the graph of \(u\) gives for its multiplier at \(y\)
\[
 B_Q(y,uy)=B_Q(y,\bar u y)=\tfrac12B_Q(y,y)=Q(y).
\]

This identifies the quadratic refinement of this particular line.
The coefficient transformation and its compatibility with the
integral metric are given in
\cite[Proposition~3.19 and Corollary~3.21]{frss-nonarch-theta}. Write \(a_z\) for its coefficient of index \(z\).
Translation by the period \(z\) carries the integral abelian-part
section at index zero to that at index \(z\), with multiplier
of order \(Q(z)\). In the integral frame this gives
\[
                 \operatorname{ord}(a_z)\ge Q(z).
\]

Write the original torsion lift as \(t\), with \([m^2]t\) equal
to the period \(y_t\) on the uniformizing group. Character orders
are additive under multiplication. At \([m]t\) the character
\(z\) contributes \(B_Q(y_t,z)/m\). The symmetric torsion frame
contributes \(Q(y_t)/m^2\): its square has valuation
\(B_Q(y_t,y_t)/m^2\) by the symmetric biextension pairing after clearing \(m^2\),
and the rigidified cubical isomorphism \([m]^*L\simeq L^{m^2}\) is integral. Thus
the term of index \(z\) has order bounded below by
\[
 Q(z)+\frac{B_Q(y_t,z)}m+\frac{Q(y_t)}{m^2}
       =\frac{Q(y_t+mz)}{m^2}.
\]
All orders are measured in the uniformizer of the same split chart.
The terms of the convergent theta expansion lie in the closed
valuation ball specified by their common lower bound. Their sum
lies in that ball as well, even when some terms cancel.

For a spanning tuple of torsion labels the classes \(y_t\) span the
degeneration lattice modulo \(R^2\). Indeed the period map from
geometric \(m^2\)-torsion to \(Y/m^2Y\) sends a lift to its
\(m^2\)-multiple modulo periods. It is surjective: each period
has an \(m^2\)-division point in the semiabelian uniformizing
group over the algebraic closure of the fraction field. It
commutes with the \(\mathcal O_E\)-action. The resulting surjections, with \(A\) the degenerating abelian
variety, are \(\mathcal O_E\)-equivariant:
\[
 A[m^2](\overline K)\twoheadrightarrow Y/m^2Y\twoheadrightarrow Y/mY .
\]
Thus their images generate \(Y/mY\). Replacing each image by a minimizing
representative preserves its class in this quotient. Since the lattice is a finite module over
\(\mathcal O_E\), these representatives span after tensoring
with \(E\) by Nakayama at a prime above \(R\).
Write \(v_b\) for these representatives. On balanced integral
symmetric bilinear forms \(C\) use the integral linear functional
\[
 \ell(C)=\sum_b C(uv_b,v_b)
        =\tfrac12\sum_b C(v_b,v_b),\qquad u+\bar u=1 .
\]

The equality follows from symmetry and
\(C(ux,y)=C(x,\bar u y)\). If \(C\ge0\) and \(\ell(C)=0\),
all \(v_b\) lie in the radical, which is \(\mathcal O_E\)-stable.
Their \(E\)-span is the whole rational lattice, so \(C=0\).
Thus \(\ell\) is one of the integral functionals allowed in the
infimum defining \(\Phi\). Evaluating the functional on \(B_Q\) and using the preceding
termwise bound gives
\[
 \Phi\le\ell(B_Q)=\sum_b Q(v_b)
       \le R^4\operatorname{ord}\!\left(\prod_b A(\gamma x_{l,b})\right).
\] The integer coefficients and the
ultrametric inequality give that same lower bound for the
whole factor. Multiplication over \(\Gamma_R\) gives
\(2i_\eta\Phi\). Descent uses the normalized valuation on the
coefficient field: if the extension ramification is \(e\),
both the raw order and the boundary form computed using the
extension's uniformizer acquire a factor \(e\).

If \(V\subset V'\) are the two DVRs, the numerical descent is
\[
  {\rm ord}_{V'}(a)=e\,{\rm ord}_{V}(a),\qquad {\rm ord}_{V}(a)=e^{-1}{\rm ord}_{V'}(a)\qquad(a\in V,\ a\ne0).
\]
In particular \(v_\lambda(\eta)\ge2i_\eta\phi_\lambda\)
in the relative uniformizer units of the finite light
comparison. A zero value has infinite order. The section
is integral at all algebraic tests and therefore belongs
to the integral cusp lattice once cuspidality is established below.
\label{cor:selected-auxiliary-cusp-4}

\proofstep{The complex norm on both halves.}
For a good complex embedding take a unit frame. With \(d\mu\) the Haar probability measure of the torus, the chosen unit frame normalizes the Appell--Humbert section \(\vartheta\) by
\[
  \int\|\vartheta(z)\|^2\,d\mu(z)=1.
\]
The absolute value of the normalized theta null value is the norm
of this section at the torsion point represented by the period
divided by \(R^2\). Multiplying these period representatives
by \(R^2\) produces a spanning family modulo \(R\) for every
tuple in \eqref{eq:selected-eta-product}. In this chart the
polarization has semicharacter \((-1)^{B(k,uk)}\). Indeed the
rigidified theta line is obtained by pulling the Poincar\'e line
back along \(z\mapsto(z,uz)\) using the polarization. Its sign
on a period is the pairing of the two period components modulo two.
The two axis rigidifications determine this sign, as in
Lemma~\ref{lem:placed-row-theta}.

Take \(g=fr\), \(M=g_0g_1\), \(R_0^2=90(g+D_F)\), and
\(\rho=\operatorname{rd}(E)\). The boundary height used on the
selected page is the function \(P\) of Lemma~\ref{lem:eta-period-decay}
at the threshold
\[
 h=\sqrt{f R_0^2\rho^{21}/M}.
\]

The page has \(D_F,m\ge1\), \(fr\ge2\), \(m^4\le f\),
\(\rho\ge\max(16,f^2)\), and
\begin{equation}\label{eq:selected-torus-level}
 M\ge R_0^2\rho^{28} D_F r .
\end{equation}

To verify the level bound on the chosen lift, use the same size
parameter \(\mathcal P=C_2z^{A_2}\) as in the proof of
Lemma~\ref{lem:short-weight-curvature}. The construction takes
\(\mathcal P\ge90(fr+1)\) and \(\mathcal P\ge r\); its prime
inequality is \(\rho^{28}\mathcal P^3 D_F^2\le\min(g_0,g_1)\).
Since \(R_0^2\le\mathcal P D_F\) and \(r\le\mathcal P^2\), we obtain
\[
 R_0^2\rho^{28}D_F r\le\rho^{28}\mathcal P^3D_F^2
                    \le\min(g_0,g_1)\le M .
\]

At the nonterminal signature the smaller block has size
\(m-j\le m\). Lemma~\ref{lem:joint-theta-parameters} now applies
both to the light-family estimate in
Lemma~\ref{lem:scaled-light-metric} and to the decay estimate at
this same threshold. In particular \(M\ge\rho^2\) and
\(M^{-1/2}\le h\le\Omega^{-1}\) as required by
Lemma~\ref{lem:eta-period-decay}.
\label{cor:selected-auxiliary-cusp-5}

Write \(A_0=R^4rg\ge1\). For \(P>1\), concavity of the logarithm
gives \(\log P\le\log A_0+P/A_0\). For \(0\le P\le1\) the same upper
bound holds for \(\log\max(1,P)=0\). Since \(\pi\ge2\),
\eqref{eq:eta-period-decay} yields, for every tuple in
\eqref{eq:selected-eta-product},
\[
 \prod_b|A(\gamma x_{l,b})|
 \le\exp\!\left(C_0+\frac{rg}2\log A_0
                          -\frac{P}{2R^4}\right).
\]
The triangular change to a unit frame preserves this boundary
function, and the flag metric gives exactly the norm of
the value in a unit frame.

The Lie factor is one at a scalar
weight. Using \(|J|\le R^{2fr^2}\) and
\(0\le m_l\le B\le B_1=\max(1,M_c)R^{8fr^2}\) in the triangle
inequality, we get
\begin{align*}
   \log\|\eta\|&\le i_\eta(p_{16}-P),\\
   p_{16}&=2R^4(C_0+\log B_1+2fr^2\log R)+A_0\log A_0.
\end{align*}
The bound applies on each good embedding, and on the others
the norms come from conjugation.

To see cuspidality directly, fix a rational boundary component and
a primitive period in its isotropic lattice. In a Witt basis one
may approach this component with the real split diagonal acting
by \(e^{-t}\) on the isotropic space and by \(e^t\) on its dual.
The unipotent radical fixes that isotropic space pointwise.
For fixed Levi coordinate and uniformly over the compact unipotent
period of integration, the first minimum therefore satisfies
\[
 \lambda_1(t)\le C e^{-t}<h,\qquad
 P(t)\ge\frac{2f}{\rho^2 C^2}e^{2t}\qquad(t\gg1).
\] The change from
the metric frame to the holomorphic Fourier frame has at most
exponential growth in \(t\): the determinant characters of the
scalar weight are fixed powers of the diagonal entries. Hence
\[
 \sup_{\rm period}|\eta|_{\rm Fourier\ frame}
      \le C_1\exp(C_2t-C_3e^{2t})\longrightarrow0
      \qquad(C_3>0).
\]

The coefficient of index zero in the normal Fourier variables is
their period average. It is independent of the imaginary
translation in those variables by holomorphy and periodicity,
so taking this limit makes it zero. This holds at each rational
boundary component. A component with no rational boundary has
no constant-term condition. Together with the finite orders this
puts \(\eta\) in the integral cusp lattice.
\label{cor:selected-auxiliary-cusp-6}

\proofstep{The size bound in the selected parameters.}
For the explicit size, denote the run's logarithmic scale by
\(L_0\) and put \(Z=frL_0\). The page has \(L_0\ge2\),
\(\log r\le L_0\), \(\log f,\log\operatorname{rd}(E)\le 2L_0\),
\(\log M\le70L_0\), \(m\le r\), and \(\log M_c\le41Z^2\).
For \(R=5\), \(r\ge2\), and \(g=fr\), they give
\[
 C_0\le fr^2(8+39L_0),\qquad
 \log A_0=4\log5+\log f+2\log r\le7+4L_0.
\]

Indeed \(\log g\le3L_0\), \(2f\le g\), and
\(\log(2\sqrt{2f})\le2+L_0\). As \(\log5<2\), it follows that
\[
 p_{16}\le1250[41Z^2+fr^2(28+39L_0)]
             +625fr^2(7+4L_0)
        \le 200000 Z^2 .
\]

For the last inequality use \(fr^2L_0^2\le Z^2\) and
\(L_0\ge2\). This supplies both the constant \(P_2\) in the
decay estimate and the constant \(P_\eta\) in the mean
log supremum divided by \(i_\eta\). The same scalar form
supplies the finite order comparison. It is selected before
the small short basis on the incidence field.
\label{cor:selected-auxiliary-cusp-7}
\end{proof}
